%! Author = elbekbakiev
%! Date = 03.01.25

\section{Abstract}
\paragraph{}
    Machine learning is actively used in processing and restoring audio signals, particularly to improve their quality and
    repair damaged audio parts. This study addresses the problem of restoring lost segments of speech
    audio signals that may contain critical information.

    To solve this task, a Python framework will be developed and utilized; NumPy and SciPy for signal analysis and
    transformation; Librosa for extracting and processing audio signal characteristics; Kaldi (or TensorFlow/PyTorch)
    for developing and training machine learning models capable of predicting missing fragments.

    The data for the study will be sourced from open-access materials. For the data processing, gaps and
    various types of noise will be artificially introduced, creating a training set for the models.
    The primary goal of the work is to develop an effective approach to restoring lost data, which is crucial
    for improving the quality of voice communication and other applications.


{\color{red}Abstract could be modified and adjusted according to the results of the work at the last step.}
